19. Tiêu chuẩn Tuyển sinh \& Quy trình Nộp hồ sơ
Chương trình MS-DSAI tìm kiếm các ứng viên có năng lực định lượng xuất sắc và nền tảng lập trình vững chắc. Việc xét tuyển được thực hiện theo hình thức cuốn chiếu (rolling basis) cho cả hai kỳ nhập học Mùa Thu và Mùa Xuân.
Yêu cầu Tiền đề Học thuật
Bằng Cử nhân Khoa học Máy tính (CS), Khoa học Dữ liệu, Toán học, Thống kê,
Vật lý, hoặc Kỹ thuật.
GPA Đại học Tối thiểu: 3.20 trên thang điểm 4.00.
Các môn học tiên quyết: Giải tích III, Đại số Tuyến tính,
Lý thuyết Xác suất, và Cấu trúc Dữ liệu bằng C++/Python.
Danh mục Hồ sơ Ứng tuyển
Bảng điểm Đại học Chính thức.
3 Thư Giới thiệu (Học thuật/Chuyên môn).
Bài luận Mục đích Học tập (1.000 từ).
Bài thi GRE General (Không bắt buộc, khuyến nghị
điểm Định lượng > 165).
TOEFL iBT (Tối thiểu 100) hoặc IELTS (Tối thiểu 7.5) đối với người không phải
là người bản ngữ tiếng Anh.
20. Học phí, Lệ phí \& Hỗ trợ Tài chính
Khoản chi phí
Chi phí trên mỗi đơn vị
Tổng chi phí toàn chương trình (36 Tín chỉ)
Học phí Sau đại học
\$1.350 mỗi tín chỉ
\$48.600
Phí Siêu máy tính \& Hạ tầng Phòng thí nghiệm
\$450 mỗi học kỳ
\$1.800
Phí Dịch vụ \& Sức khỏe Sinh viên Sau đại học
\$350 mỗi học kỳ
\$1.400
Tổng Học phí \& Lệ phí Ước tính
—
\$51.800 USD
Hỗ trợ Tài chính Sau đại học \& Học bổng Khuyến học
Trợ lý Nghiên cứu (RA)
Đài thọ 100\% học phí + trợ cấp hàng tháng \$2.800.
Được xét duyệt cạnh tranh dựa trên sự phù hợp về hướng nghiên cứu với
các dự án tài trợ của giảng viên. Yêu cầu 20 giờ/tuần làm việc tại phòng thí nghiệm.
Trợ giảng (TA)
Đài thọ từ 50\% đến 100\% học phí + trợ cấp \$2.400/tháng.
Các TA hỗ trợ giảng viên trong việc chấm điểm, hướng dẫn các buổi
thực hành phòng lab, và chủ trì các giờ tư vấn sinh viên.
Thạc sĩ Khoa học Dữ liệu \& Trí tuệ Nhân tạo — Sổ tay Học thuật
Trang 11 / 12

21. Định hướng Nghề nghiệp \& Việc làm Sau Đại học
Sinh viên tốt nghiệp chương trình MS-DSAI được tuyển dụng bởi các công ty công nghệ hàng đầu thế giới, các phòng thí nghiệm nghiên cứu, các tổ chức tài chính, và các công ty khởi nghiệp sáng tạo có tốc độ phát triển cao.
\$148k
MỨC LƯƠNG KHỞI ĐIỂM TRUNG BÌNH
98\%
CÓ VIỆC LÀM TRONG VÒNG 6 THÁNG
45+
ĐỐI TÁC TUYỂN DỤNG DOANH NGHIỆP
12\%
TIẾP TỤC HỌC CHƯƠNG TRÌNH TIẾN SĨ
Các Vị trí Tuyển dụng Hàng đầu \& Đối tác Doanh nghiệp Chính
Các Chức danh Công việc Mục tiêu Phổ biến
Senior Machine Learning Engineer (Kỹ sư Học máy Cao cấp)
AI Systems Architect (Kiến trúc sư Hệ thống AI)
Principal Data Scientist (Chuyên gia Khoa học Dữ liệu Chính)
MLOps Infrastructure Engineer (Kỹ sư Hạ tầng MLOps)
AI Regulatory \& Governance Specialist (Chuyên gia Quản trị \& Pháp lý AI)
Computer Vision / NLP Research Scientist (Nhà nghiên cứu Thị giác Máy tính / Xử lý Ngôn ngữ Tự nhiên)
Mạng lưới Đối tác Tuyển dụng
Google • Microsoft • Amazon Web Services • NVIDIA • Meta AI
• OpenAI • Snowflake • Databricks • Goldman Sachs •
JPMorgan Chase • Mayo Clinic Digital Health • Anthropic
22. Danh mục Kiểm tra Điều kiện Tốt nghiệp
Các Bước Hoàn thành Bằng cấp Cuối cùng:
1. Hoàn thành toàn bộ 36 tín chỉ bắt buộc với GPA tích lũy ≥ 3.00.
2. Nộp và bảo vệ thành công trước công chúng Luận văn / Đề án Tốt nghiệp (Capstone).
3. Hoàn tất xác nhận hoàn thành nghĩa vụ (clearance) từ Thư viện Trường và Cơ sở Siêu máy tính.
4. Nộp Đơn đăng ký Tốt nghiệp với Phòng Đào tạo (Office of the Registrar).
23. Thông tin Liên hệ Đơn vị
Để biết thêm thông tin chi tiết về chương trình học thuật, yêu cầu tuyển sinh hoặc hợp tác nghiên cứu, vui lòng liên hệ văn phòng chương trình:
Phòng Tuyển sinh Sau Đại học
Trường Máy tính, Viện Công nghệ Apex
100 Innovation Boulevard, Tech District
Email: admissions-dsai@apex.edu
Phòng Thường trực Học vụ
Điện thoại: +1 (555) 019-2831
Giờ làm việc: Thứ Hai – Thứ Sáu, 8:30 sáng – 5:00 chiều EST
Web: https://dsai.apex.edu/handbook
Thạc sĩ Khoa học Dữ liệu \& Trí tuệ Nhân tạo — Sổ tay Học thuật
Trang 12 / 12
